\section{Recover samples from residuals}
Each inflated row starts with a filter selector from zero through four. The remaining
bytes are residuals: differences from predicted sample values, reduced modulo 256.
Reconstruction adds the same predictor back and keeps the low eight bits. Let $a$ be
the already reconstructed sample three bytes to the left, $b$ the sample above, and
$c$ the upper-left sample. Use zero when a neighbor does not exist. The three-byte
spacing follows this chapter's RGB8 pixel format.

\Needspace{10\baselineskip}
\begin{longtable}{p{18mm}p{28mm}p{104mm}}
  \toprule Selector & Name    & Predictor added to the residual                     \\\midrule
  0                 & None    & Zero                                                \\
  1                 & Sub     & $a$                                                 \\
  2                 & Up      & $b$                                                 \\
  3                 & Average & $\lfloor(a+b)/2\rfloor$                             \\
  4                 & Paeth   & One of $a$, $b$, $c$, selected by the distance rule \\\bottomrule
\end{longtable}
The predictor uses reconstructed neighbors, not their encoded residuals. Exclude the
filter selector itself from sample indexing. If rows remain in the inflated buffer,
the previous sample row begins one byte after its selector, and row starts are seven
bytes apart in the worked image. If samples move into a compact RGB buffer instead,
row starts are six bytes apart. Choose one representation before deriving addresses.

For Average, widen the unsigned byte values before adding: $250+240=490$, so the
predictor is 245. Wrapping the sum to a byte first would produce 117. Only the final
residual-plus-predictor addition wraps to eight bits. This distinction is easy to lose
when replacing scalar code with byte-lane operations.
